\manualchapter{Control reference}{sec:controls}
The left column contains shared controls. The other four columns use identical
operator layouts and natural algorithm order 1, 2, 3, 4 (hardware register
slots 0, 2, 1, 3). The algorithm display shows modulation connections in teal
and audible carriers in gold. Feedback always belongs to operator 1.

All numeric chip controls round to the nearest integer after adding CV,
then clamp. A +8\,V input adds the full positive range; negative CV subtracts.
Unpatched CV is zero. TUNE and output LEVEL are continuous. Every parameter
is saved by Rack, including discrete choices; no custom JSON is required.

\begin{longtable}{@{}p{.16\textwidth}p{.25\textwidth}p{.50\textwidth}@{}}
\toprule Control & Range; default & Meaning and CV \\\midrule\endhead
TUNE & $-4$--4 octaves; 0 & Adds to V/OCT; 0\,V is C4. Pitch input is 1\,V/octave. \\
ALG / FB & 0--7; 0 / 0 & Eight algorithms / operator-1 feedback; dedicated CV each. \\
LEVEL & 0--1; 0.8 & Output gain; dedicated CV adds 1 per 8\,V. \\
LFO & 0--255; 0 & Native frequency register, faster at larger values; CV. \\
AMD / PMD & 0--127; 0 & Separate LFO amplitude/pitch depths; dedicated CV each. \\
AMS / PMS & 0--3 / 0--7; 0 & Channel modulation sensitivities; no CV. \\
WAVE & 0--3; saw & Saw, square, triangle, noise LFO; no CV. \\
NOISE / RATE & Off / 0--31; off / 16 & Replace operator 4's sine with noise. RATE has CV; larger values are faster. \\
ROUTE & Left, Right, Both; Both & Hardware channel routing; no CV. \\
\bottomrule
\end{longtable}

\clearpage
\begin{longtable}{@{}p{.16\textwidth}p{.25\textwidth}p{.50\textwidth}@{}}
\toprule Per operator & Range; default & Meaning and CV \\\midrule\endhead
AR & 0--31; 31 & Attack rate; 0 prevents attack from silence. CV. \\
D1 / SR & 0--31; 0 & First decay / sustain decay rates. Zero holds that stage. CV each. \\
SL & 0--15; 15 & Sustain level, inverted to chip attenuation (15 minus knob). CV. \\
RR & 0--15; 15 & Release rate. Zero is slowest, not an infinite hold. CV. \\
VOL & 0--127; 100,100,100,127 & Level, inverted to chip total attenuation (127 minus knob), in 0.75\,dB steps. CV. \\
MUL & 0--15; 1 & Frequency multiplier; register 0 means $0.5\times$, others $1$--$15\times$. CV. \\
KS & 0--3; 0 & Key-rate scaling, increasing rate dependence on pitch; no CV. \\
DT1 & 0--7; 0 & Native fine detune code: 0/4 zero, 1--3 positive, 5--7 negative. No CV. \\
DT2 & 0--3; 0 & Coarse detune: 0, 600, 781, 950 cents. No CV. \\
AM & 0--1; off & Enable LFO amplitude modulation for this operator; no CV. \\
\bottomrule
\end{longtable}

The native sustain attenuation steps are 3\,dB; chip code 15 jumps to 93\,dB.
Thus panel SL 0 selects that lowest threshold, while SL 15 selects 0\,dB.
Rates are register values, not seconds. Higher rates are faster, and KS/pitch
can change envelope timing. AM requires nonzero AMD and AMS as well as an
operator's AM switch. PM requires both PMD and PMS. Setting depths alone
therefore need not change the sound.

GATE rises at 2\,V and falls at or below 0.01\,V. RETRIG uses the same
hysteresis. With gate held high, one retrigger edge schedules key-off then
key-on. A held trigger does not repeat. Simultaneous gate rise and retrigger
key on once; gate fall wins over retrigger; retrigger alone cannot start a note.
